% !TEX root = ctfp-print.tex

\input{half-title}

\frontmatter
\tableofcontents

\chapter*{A note from the editor}
\addcontentsline{toc}{chapter}{A note from the editor}
% !TEX root = ctfp-print.tex

\lettrine[lhang=0.17]{T}{his is the Scala edition} of \emph{Category Theory for Programmers}.
It's been a tremendous success, making Bartosz Milewski's blog post series available as a nicely 
typeset \acronym{PDF}, as well as a hardcover book. There have been numerous contributions made
to improve the book, by fixing typos and errors, as well as translating the code snippets into
other programming languages.

I am thrilled to present this edition of the book, containing the original Haskell code, followed by
its Scala counterpart. The Scala code snippets were generously provided by 
\urlref{https://github.com/typelevel/CT_from_Programmers.scala}{Typelevel} contributors, slightly 
modified to suit the format of this book.

To support code snippets in multiple languages, I am using a \LaTeX{} macro to load the code snippets
from external files. This allows easily extending the book with other languages, while leaving the 
original text intact. Which is why you should mentally append the words ``and Scala'' whenever you see
``in Haskell'' in the text.

The code is laid out in the following manner: the original Haskell code, followed by Scala code.
To distinguish between them, the code snippets are braced from the left with a vertical bar, in the primary
color of the language's logo, \raisebox{-.2mm}{\includegraphics[height=.3cm]{fig/icons/haskell.png}}, 
and \raisebox{-.2mm}{\includegraphics[height=.3cm]{fig/icons/scala.png}} respectively, e.g.:

\srcsnippet{content/3.6/code/haskell/snippet03.hs}{blue}{haskell}
\unskip
\srcsnippet{content/3.6/code/scala/snippet03.scala}{red}{scala}
\NoIndentAfterThis

In addition, some Scala snippets make use of the 
\urlref{https://github.com/non/kind-projector}{Kind Projector} compiler plugin, to support nicer syntax for partially-applied types.
\chapter*{Preface}
\addcontentsline{toc}{chapter}{Preface}
\label{Preface}
\subfile{content/0.0/Preface}

\mainmatter

\part*{Part One}
\addcontentsline{toc}{part}{Part One}

\chapter{Category: The Essence of Composition}\label{category-the-essence-of-composition}
\subfile{content/1.1/"Category - The Essence of Composition"}

\chapter{Types and Functions}\label{types-and-functions}
\subfile{content/1.2/"Types and Functions"}

\chapter{Categories Great and Small}\label{categories-great-and-small}
\subfile{content/1.3/"Categories Great and Small"}

\chapter{Kleisli Categories}\label{kleisli-categories}
\subfile{content/1.4/"Kleisli Categories"}

\chapter{Products and Coproducts}\label{products-and-coproducts}
\subfile{content/1.5/"Products and Coproducts"}

\chapter{Simple Algebraic Data Types}\label{simple-algebraic-data-types}
\subfile{content/1.6/"Simple Algebraic Data Types"}

\chapter{Functors}\label{functors}
\subfile{content/1.7/Functors}

\chapter{Functoriality}\label{functoriality}
\subfile{content/1.8/Functoriality}

\chapter{Function Types}\label{function-types}
\subfile{content/1.9/"Function Types"}

\chapter{Natural Transformations}\label{natural-transformations}
\subfile{content/1.10/"Natural Transformations"}

\part*{Part Two}
\addcontentsline{toc}{part}{Part Two}

\chapter{Declarative Programming}\label{declarative-programming}
\subfile{content/2.1/"Declarative Programming"}

\chapter{Limits and Colimits}\label{limits-and-colimits}
\subfile{content/2.2/"Limits and Colimits"}

\chapter{Free Monoids}\label{free-monoids}
\subfile{content/2.3/"Free Monoids"}

\chapter{Representable Functors}\label{representable-functors}
\subfile{content/2.4/"Representable Functors"}

\chapter{The Yoneda Lemma}\label{the-yoneda-lemma}
\subfile{content/2.5/"The Yoneda Lemma"}

\chapter{Yoneda Embedding}\label{yoneda-embedding}
\subfile{content/2.6/"Yoneda Embedding"}

\part*{Part Three}
\addcontentsline{toc}{part}{Part Three}

\chapter{It's All About Morphisms}\label{all-about-morphisms}
\subfile{content/3.1/"It's All About Morphisms"}

\chapter{Adjunctions}\label{adjunctions}
\subfile{content/3.2/Adjunctions}

\chapter{Free/Forgetful Adjunctions}\label{free-forgetful-adjunctions}
\subfile{content/3.3/"Free-Forgetful Adjunctions"}

\chapter{Monads: Programmer's Definition}\label{monads-programmers-definition}
\subfile{content/3.4/"Monads - Programmer's Definition"}

\chapter{Monads and Effects}\label{monads-and-effects}
\subfile{content/3.5/"Monads and Effects"}

\chapter{Monads Categorically}\label{monads-categorically}
\subfile{content/3.6/"Monads Categorically"}

\chapter{Comonads}\label{comonads}
\subfile{content/3.7/Comonads}

\chapter{F-Algebras}\label{f-algebras}
\subfile{content/3.8/F-Algebras}

\chapter{Algebras for Monads}\label{algebras-for-monads}
\subfile{content/3.9/"Algebras for Monads"}

\chapter{Ends and Coends}\label{ends-and-coends}
\subfile{content/3.10/"Ends and Coends"}

\chapter{Kan Extensions}\label{kan-extensions}
\subfile{content/3.11/"Kan Extensions"}

\chapter{Enriched Categories}\label{enriched-categories}
\subfile{content/3.12/"Enriched Categories"}

\chapter{Topoi}\label{topoi}
\subfile{content/3.13/Topoi}

\chapter{Lawvere Theories}\label{lawvere-theories}
\subfile{content/3.14/"Lawvere Theories"}

\chapter{Monads, Monoids, and Categories}\label{monads-monoids-categories}
\subfile{content/3.15/"Monads, Monoids, and Categories"}

\backmatter

\appendix
\addcontentsline{toc}{part}{Appendices}
\input{index}

\makeatletter\@openrightfalse
\chapter*{Acknowledgments}\label{acknowledgments}
\addcontentsline{toc}{chapter}{Acknowledgments}
\input{acknowledgments}

\chapter*{Colophon}\label{colophon}
\addcontentsline{toc}{chapter}{Colophon}
\lettrine[lraise=-0.03,loversize=0.08]{T}{his book} was compiled by \urlref{https://hmemcpy.com}{Igal Tabachnik}, by converting the original text by Bartosz Milewski into \LaTeX{} format, by
first scraping the original WordPress blog posts using \urlref{https://mercury.postlight.com/web-parser/}{Mercury Web Parser}
to get a clean HTML content, modifying and tweaking with \urlref{https://www.crummy.com/software/BeautifulSoup/}{Beautiful Soup},
finally, converting to LaTeX with \urlref{https://pandoc.org/}{Pandoc}.

The typefaces are Linux Libertine for body text and Linux Biolinum for headings, both by Philipp H. Poll. Typewriter face is Inconsolata
created by Raph Levien and supplemented by Dimosthenis Kaponis and Takashi Tanigawa in the form of Inconsolata \acronym{LGC}. The cover page
typeface is Alegreya, designed by Juan Pablo del Peral.

Original book layout design and typography are done by Andres Raba. Syntax highlighting is using ``GitHub'' style for Pygments by
\urlref{https://github.com/hugomaiavieira/pygments-style-github}{Hugo Maia Vieira}. Scala code translation was done by 
\urlref{https://github.com/typelevel/CT_from_Programmers.scala}{Typelevel} contributors.

\chapter*{Copyleft notice}\label{copyleft}
\addcontentsline{toc}{chapter}{Copyleft notice}
\input{free-software}
\@openrighttrue\makeatother
\afterpage{\blankpage}

\input{postamble}